\documentclass[11pt]{article}

\usepackage[letterpaper,margin=0.82in]{geometry}
\usepackage{amsmath,amssymb,mathtools}
\usepackage{microtype}
\usepackage{xcolor}
\usepackage[most]{tcolorbox}
\usepackage{fancyhdr}
\usepackage{enumitem}

\definecolor{navy}{HTML}{17365D}
\definecolor{blue}{HTML}{2D6CA2}
\definecolor{pale}{HTML}{F3F7FB}
\definecolor{ink}{HTML}{18212B}
\definecolor{muted}{HTML}{637083}
\definecolor{green}{HTML}{247A4A}

\color{ink}
\setlength{\parindent}{0pt}
\setlength{\parskip}{0.55em}
\allowdisplaybreaks

\pagestyle{fancy}
\fancyhf{}
\lhead{\footnotesize\color{muted}MathForge technique-combination pilot}
\rhead{\footnotesize\color{muted}Algebra}
\cfoot{\footnotesize\color{muted}\thepage}
\renewcommand{\headrulewidth}{0.3pt}

\newtcolorbox{problembox}{
  enhanced,
  breakable,
  colback=pale,
  colframe=blue,
  boxrule=0.9pt,
  arc=2.5mm,
  left=5mm,right=5mm,top=4mm,bottom=4mm,
  title={\bfseries Problem},
  coltitle=white,
  colbacktitle=blue,
  attach boxed title to top left={xshift=4mm,yshift=-2mm},
  boxed title style={arc=1.5mm,boxrule=0pt}
}

\newtcolorbox{resultbox}{
  colback=green!6,
  colframe=green,
  boxrule=0.8pt,
  arc=2mm,
  left=4mm,right=4mm,top=3mm,bottom=3mm
}

\newcommand{\solutionhead}[1]{%
  \par\medskip{\large\bfseries\color{ink}#1}\par\smallskip
}

\begin{document}

\begin{center}
  {\small\bfseries\color{blue}COMBINED-TECHNIQUE COMPETITION PROBLEM}\par
  \vspace{0.45em}
  {\LARGE\bfseries\color{navy}Lagrange Interpolation and Coefficient Comparison}\par
  \vspace{0.45em}
  {\color{muted}Target level: late AIME / introductory olympiad}\par
\end{center}

\vspace{1em}

\begin{problembox}
Let $t$ be a positive integer different from each of $1,2,3,4,6$, so that the
six numbers $1,2,3,4,6,t$ are distinct. For each of these six numbers $v$, let
$D(v)$ denote the product of the five differences $(v-u)$ taken over the other
five numbers $u$. Suppose that
\[
  \frac{1^{7}}{D(1)}+
  \frac{2^{7}}{D(2)}+
  \frac{3^{7}}{D(3)}+
  \frac{4^{7}}{D(4)}+
  \frac{6^{7}}{D(6)}+
  \frac{t^{7}}{D(t)}=497.
\]
Find $t$.
\end{problembox}

\vspace{1.2em}

\begin{center}
\begin{minipage}{0.82\textwidth}
\small\color{muted}
\textbf{Review note.} The problem appears alone on this page so it can be read
without spoilers. A complete solution begins on the next page.
\end{minipage}
\end{center}

\vfill

\begin{center}
  \small\color{muted}
  Claude panel: $3/3$ answer agreement, well-posedness accepted, difficulty $6.5/10$,
  elegance $5/5$.
\end{center}

\newpage

\begin{center}
  {\Large\bfseries\color{navy}Solution}\par
\end{center}

\begingroup
\small
\setlength{\parskip}{0.28em}

Write the six nodes as
\[
  \{x_0,x_1,\ldots,x_5\}=\{1,2,3,4,6,t\},
  \qquad
  w_i=\frac{1}{\prod_{j\ne i}(x_i-x_j)}=\frac{1}{D(x_i)}.
\]
The Lagrange basis polynomial
\[
  L_i(x)=w_i\prod_{j\ne i}(x-x_j)
\]
has degree $5$ and leading coefficient $w_i$.

\solutionhead{A generating identity}

Consider
\[
  G(z)=\sum_i w_i\prod_{k\ne i}(1-x_kz),
\]
a polynomial of degree at most $5$. Since every node is nonzero, evaluate at
$z=1/x_\ell$. Every term with $i\ne\ell$ vanishes, so
\begin{align*}
  G(1/x_\ell)
    &=w_\ell\prod_{k\ne\ell}\left(1-\frac{x_k}{x_\ell}\right)\\
    &=w_\ell\frac{\prod_{k\ne\ell}(x_\ell-x_k)}{x_\ell^5}
      =\frac{1}{x_\ell^5}.
\end{align*}
Thus $G(z)-z^5$ has degree at most $5$ and vanishes at the six distinct points
$1/x_i$. Therefore $G(z)=z^5$. Dividing by $\prod_i(1-x_iz)$ gives
\[
  \sum_i\frac{w_i}{1-x_iz}
  =\frac{z^5}{\prod_i(1-x_iz)}
  =z^5\sum_{k\ge0}h_kz^k,
\]
where $h_k$ is the complete homogeneous symmetric polynomial of degree $k$ in
the six nodes.

\solutionhead{Comparing coefficients}

The coefficient of $z^m$ on the left is $\sum_i w_ix_i^m$, while the
coefficient on the right is $h_{m-5}$. In particular, when $m=7$,
\[
  \sum_i\frac{x_i^7}{D(x_i)}=\sum_iw_ix_i^7=h_2.
\]
The left side is exactly the sum in the problem.

For six variables,
\[
  h_2=\frac{e_1^2+p_2}{2}.
\]
Here
\[
  e_1=1+2+3+4+6+t=16+t,
  \qquad
  p_2=1+4+9+16+36+t^2=66+t^2.
\]
Consequently,
\begin{align*}
  h_2
    &=\frac{(16+t)^2+66+t^2}{2}\\
    &=t^2+16t+161.
\end{align*}

\solutionhead{Solving for the node}

Setting $h_2=497$ gives
\[
  t^2+16t-336=(t-12)(t+28)=0.
\]
The only positive solution is $t=12$, which is distinct from
$1,2,3,4,6$.

\begin{resultbox}
\centering
\large Therefore, the answer is \(\boxed{12}\).
\end{resultbox}

\endgroup

\end{document}
